% =============================================================================
%  ABSTRACT
% =============================================================================
\newpage
\begin{resumo}[Abstract]
\SingleSpacing

This work presents a systematic and in-depth investigation into the application of Digital Twins in dentistry, with emphasis on the development of the SUS-Twin framework for periodontal biomechanical simulation within the Brazilian Unified Health System (SUS). The book is structured in four complementary parts: (i) fundamentals of digital twins, covering definitions, architectures, and medical applications; (ii) digital dentistry, exploring orthodontics, implantology, endodontics, prosthodontics, and simulation-based dental education; (iii) artificial intelligence ecosystems and open platforms, detailing the OpenCode Ecosystem, its 3D reconstruction pipelines, K-Fold methods for periodontal prediction, and IoT gateways for telemetry; and (iv) social impact and future perspectives, analyzing equity in SUS, ethics, privacy, regulatory frameworks, and biomechanical challenges.

The methodology combines systematic literature review (20+ publications 2023-2026), viscoelastic mathematical modeling of the periodontal ligament (Prony series), K-Fold cross-validation with RMSE $< 0.15$ mm, and zero-knowledge proofs (ZKP) for cryptographic health record auditing. The SUS-Twin framework was implemented in Python with a suite of 11 TDD tests (100\% approval rate) and an interactive Streamlit dashboard. Social Return on Investment (SROI) analysis indicates an average ratio of $3.5:1$, demonstrating the economic-social viability of the technology for the Brazilian public health system.

\textbf{Keywords}: Digital Twins. Digital Dentistry. Artificial Intelligence. SUS-Twin. OpenCode Ecosystem. Periodontal Biomechanics. Health Equity.
\end{resumo}
